% !TeX root = main.tex
% PAGINE ORIGINALI 511-520 — DA COMPILARE
% Incollare qui i quiz del blocco e alla fine le rispettive\registraSoluzione{numero}{lettera}.
% Questo file vuoto non aggiunge pagine al PDF.

% !TeX root = main.tex
% Pagine 511-520: tabella risposte simulazione 2; simulazioni 3 e 4 del giugno 2012.

% Pag. originale 512/521
\iffalse Pag. originale 512/521 \fi
\begin{quiz}
Il dominio della funzione $f(x)=\sqrt{\log_{1/2}(x-1)+2}$ \`e:
\begin{enumerate}
\item $[5,+\infty)$
\item $[1,5]$
\item $(1,5]$
\item $\mathbb R\setminus\{1\}$
\item $(5,+\infty)$
\end{enumerate}
\end{quiz}
% SOLUZIONE: C
% SEZIONE_AM1: 1.4
% SPIEGAZIONE: Serve $x>1$ e $\log_{1/2}(x-1)\ge-2$. Essendo la base minore di $1$, si ottiene $x-1\le4$, quindi $1<x\le5$.

% Pag. originale 512/521
\iffalse Pag. originale 512/521 \fi
\begin{quiz}
Dati gli insiemi $A=\{z\in\mathbb C:|z-i|<1\}$ e $B=\{z\in\mathbb C:\operatorname{Re}(z+i)>5\}$, \`e vero che
\begin{enumerate}
\item $A\cup B$ \`e tutto il piano complesso
\item $B\setminus A=A\setminus B$
\item $A\cap B=\varnothing$
\item $A\subseteq B$
\item $A=B$
\end{enumerate}
\end{quiz}
% SOLUZIONE: C
% SEZIONE_AM1: 1.5
% SPIEGAZIONE: In $A$ la parte reale \`e compresa tra $-1$ e $1$; in $B$ \`e maggiore di $5$. I due insiemi sono disgiunti.

% Pag. originale 512/521
\iffalse Pag. originale 512/521 \fi
\begin{quiz}
Il limite $\displaystyle\lim_{x\to-\infty}\frac{e^{2x}+5x+\cos x}{\sin x-\log|x|-x}$
\begin{enumerate}
\item vale $0$
\item vale $-\infty$
\item vale $-5$
\item vale $+\infty$
\item non esiste
\end{enumerate}
\end{quiz}
% SOLUZIONE: C
% SEZIONE_AM1: 2.2
% SPIEGAZIONE: Dividendo per $x$, numeratore e denominatore tendono rispettivamente a $5$ e $-1$; il rapporto tende a $-5$.

% Pag. originale 512/521
\iffalse Pag. originale 512/521 \fi
\begin{quiz}
Sia $f:\mathbb R\to\mathbb R$ tale che $\lim_{x\to-\infty}f(x)=-5$. Allora si pu\`o concludere che
\begin{enumerate}
\item $f$ \`e una funzione limitata
\item $\exists M>0$ tale che $f(( -\infty,-M))$ \`e un insieme limitato
\item $\inf_{\mathbb R}f=-5$
\item $f(x)<0$ per ogni $x\in\operatorname{dom}f$
\item $\forall M>0$, $f((-\infty,-M))$ \`e un insieme limitato
\end{enumerate}
\end{quiz}
% SOLUZIONE: B
% SEZIONE_AM1: 2.2
% SPIEGAZIONE: La definizione di limite finito all'infinito implica che, scegliendo $M$ sufficientemente grande, i valori $f(x)$ per $x<-M$ appartengono a un intervallo limitato intorno a $-5$.

% Pag. originale 512/521
\iffalse Pag. originale 512/521 \fi
\begin{quiz}
Sia $f(x)=\begin{cases}\dfrac{1-\cosh x}{2x^2}&\text{se }x\ne0,\\5&\text{se }x=0.\end{cases}$ Allora il limite $\lim_{x\to0}f(x)$
\begin{enumerate}
\item vale $\frac12$
\item vale $0$
\item vale $5$
\item non esiste
\item vale $-\frac14$
\end{enumerate}
\end{quiz}
% SOLUZIONE: E
% SEZIONE_AM1: 2.2
% SPIEGAZIONE: $\cosh x=1+x^2/2+o(x^2)$, perci\`o $(1-\cosh x)/(2x^2)\to-1/4$. Il valore assegnato in $0$ non conta.

% Pag. originale 512/521
\iffalse Pag. originale 512/521 \fi
\begin{quiz}
Il limite della successione $a_n=\left(\dfrac{n+1}{n-1}\right)^n$ vale
\begin{enumerate}
\item $1$
\item $e$
\item $e^{-1}$
\item $e^2$
\item $+\infty$
\end{enumerate}
\end{quiz}
% SOLUZIONE: D
% SEZIONE_AM1: 2.2
% SPIEGAZIONE: $a_n=(1+2/(n-1))^n\to e^2$.

% Pag. originale 513/521
\iffalse Pag. originale 513/521 \fi
\begin{quiz}
Se $a_n$ \`e una successione monotona, allora per $n\to\infty$
\begin{enumerate}
\item $a_n$ \`e convergente
\item $a_n$ non \`e indeterminata
\item $a_n$ \`e indeterminata
\item $a_n$ \`e divergente
\item $a_n$ non \`e convergente, n'e divergente
\end{enumerate}
\end{quiz}
% SOLUZIONE: B
% SEZIONE_AM1: 2.2
% SPIEGAZIONE: Ogni successione monotona ha limite nell'insieme dei reali estesi: pu\`o convergere oppure divergere a $\pm\infty$, ma non oscillare senza limite.

% Pag. originale 513/521
\iffalse Pag. originale 513/521 \fi
\begin{quiz}
La derivata prima della funzione $f(x)=\log\!\left(\dfrac{\sqrt{1+x^2}}{x}\right)-\dfrac{1}{2x^2}$ \`e
\begin{enumerate}
\item $f'(x)=\dfrac{1}{x^3(x^2+1)}$
\item $f'(x)=\dfrac{1}{x^2(x^2+1)}$
\item $f'(x)=-\dfrac{1}{x^3(x^2+1)}$
\item $f'(x)=-\dfrac{1}{x^3(x^2-1)}$
\item nessuna delle altre risposte \`e corretta
\end{enumerate}
\end{quiz}
% SOLUZIONE: A
% SEZIONE_AM1: 4.1
% SPIEGAZIONE: Per $x>0$, $f=\frac12\log(1+x^2)-\log x-1/(2x^2)$. Quindi $f'=x/(1+x^2)-1/x+1/x^3=1/[x^3(1+x^2)]$.

% Pag. originale 513/521
\iffalse Pag. originale 513/521 \fi
\begin{quiz}
Nell'intervallo $[-1,3]$ la funzione $f(x)=x^3+2x-3$
\begin{enumerate}
\item soddisfa le ipotesi del teorema di Rolle
\item non soddisfa le ipotesi del teorema di Lagrange
\item ha un unico punto di Lagrange $c=\sqrt{7/3}$
\item ha due punti di Lagrange $c_1=\sqrt{7/3}$ e $c_2=-\sqrt{7/3}$
\item nessuna delle altre risposte \`e corretta
\end{enumerate}
\end{quiz}
% SOLUZIONE: C
% SEZIONE_AM1: 4.2
% SPIEGAZIONE: Il coefficiente della secante \`e $(f(3)-f(-1))/4=9$; $f'(c)=3c^2+2=9$ d\`a $c=\pm\sqrt{7/3}$, ma solo quello positivo appartiene a $(-1,3)$.

% Pag. originale 513/521
\iffalse Pag. originale 513/521 \fi
\begin{quiz}
Sull'intervallo $[-2,4]$ la funzione $f(x)=\sqrt[3]{x^2}$
\begin{enumerate}
\item non ha massimo assoluto
\item non ha minimo assoluto
\item ha un punto di massimo assoluto
\item ha un punto di minimo relativo non assoluto
\item ha un punto di flesso a tangente verticale
\end{enumerate}
\end{quiz}
% SOLUZIONE: C
% SEZIONE_AM1: 4.3
% SPIEGAZIONE: La funzione $|x|^{2/3}$ \`e continua sul compatto, ha minimo assoluto $0$ in $x=0$ e massimo assoluto in $x=4$.

% Pag. originale 513/521
\iffalse Pag. originale 513/521 \fi
\begin{quiz}
Il valore del parametro $a$ per cui la funzione $f(x)=x^{2a}\log(1+2x^a)$ risulta un infinitesimo di ordine $5$ per $x\to0$ \`e
\begin{enumerate}
\item $3$
\item $5/3$
\item $5$
\item $5/2$
\item $-5/3$
\end{enumerate}
\end{quiz}
% SOLUZIONE: B
% SEZIONE_AM1: 3.4
% SPIEGAZIONE: Per $a>0$, $\log(1+2x^a)\sim2x^a$, quindi $f(x)\sim2x^{3a}$: $3a=5$ implica $a=5/3$.

% Pag. originale 513/521
\iffalse Pag. originale 513/521 \fi
\begin{quiz}
Quale delle seguenti affermazioni \`e soddisfatta dalla funzione $f(x)=(x+2)\sqrt{x}$?
\begin{enumerate}
\item $\lim_{x\to0^+}f'(x)=0$
\item $f(x)$ \`e infinita di ordine $1/2$ per $x\to+\infty$
\item $f(x)$ \`e continua e derivabile nel suo dominio
\item $f(x)$ \`e infinitesima di ordine $3/2$ per $x\to0$
\item $f(x)$ ha un punto a tangente verticale in $x=0$
\end{enumerate}
\end{quiz}
% SOLUZIONE: E
% SEZIONE_AM1: 4.1
% SPIEGAZIONE: $f(x)\sim2\sqrt{x}$ e $f'(x)=\sqrt{x}+(x+2)/(2\sqrt{x})\to+\infty$ per $x\to0^+$: tangente verticale.

% Pag. originale 514/521
\iffalse Pag. originale 514/521 \fi
\begin{quiz}
Siano $f:\mathbb R\to\mathbb R$ di classe $C^2$ e $x_0\in\mathbb R$. Affinch'e $x_0$ sia punto di massimo per $f$, la condizione $f'(x_0)=0$ e $f''(x_0)<0$ \`e:
\begin{enumerate}
\item sufficiente ma non necessaria
\item necessaria e sufficiente
\item necessaria ma non sufficiente
\item n'e necessaria, n'e sufficiente
\item necessaria se esiste anche $f'''(x_0)$
\end{enumerate}
\end{quiz}
% SOLUZIONE: A
% SEZIONE_AM1: 4.3
% SPIEGAZIONE: Il test della derivata seconda \`e sufficiente per un massimo locale stretto; non \`e necessario, come mostra $f(x)=-x^4$ in $0$.

% Pag. originale 514/521
\iffalse Pag. originale 514/521 \fi
\begin{quiz}
La parte principale della funzione $f(x)=e^{1-\cos x}-x\sin(2x)-1$ per $x\to0$ \`e
\begin{enumerate}
\item $\frac83x^4$
\item $-\frac32x^2$
\item $6x^2$
\item $\frac83x^4+o(x^4)$
\item $6x^2+o(x^2)$
\end{enumerate}
\end{quiz}
% SOLUZIONE: B
% SEZIONE_AM1: 3.4
% SPIEGAZIONE: $e^{1-\cos x}-1=x^2/2+O(x^4)$ e $x\sin(2x)=2x^2+O(x^4)$: $f(x)\sim-3x^2/2$.

% Pag. originale 514/521
\iffalse Pag. originale 514/521 \fi
\begin{quiz}
Lo sviluppo di Mac Laurin di $f(x)=\sqrt{1+x}$ di ordine $3$ \`e
\begin{enumerate}
\item $f(x)=1+\frac12x-\frac18x^2+\frac1{16}x^3$
\item $f(x)=\frac1{16}x^3+o(x^3)$
\item $f(x)=1+\frac12x+\frac12x^2+\frac12x^3+o(x^3)$
\item $f(x)=1+\frac12x-\frac18x^2+o(x^3)$
\item $f(x)=1+\frac12x-\frac18x^2+\frac1{16}x^3+o(x^3)$
\end{enumerate}
\end{quiz}
% SOLUZIONE: E
% SEZIONE_AM1: 4.2
% SPIEGAZIONE: Per il binomio generalizzato $(1+x)^{1/2}=1+x/2-x^2/8+x^3/16+o(x^3)$. L'item A riporta il solo polinomio senza resto.

% Pag. originale 514/521
\iffalse Pag. originale 514/521 \fi
\begin{quiz}
Se $f$ ha sviluppo di Mac Laurin $f(x)=1-7x^3+o(x^4)$, allora NON \`e vero che
\begin{enumerate}
\item $f(x)$ ha un flesso in $x=0$
\item $f'(x)$ ha un massimo in $x=0$
\item $f''(x)$ \`e crescente nell'intorno di $x=0$
\item $f'''(0)=-42$
\item $f^{(4)}(0)=0$
\end{enumerate}
\end{quiz}
% SOLUZIONE: C
% SEZIONE_AM1: 4.2
% SPIEGAZIONE: La risposta ufficiale \`e C: il termine cubico implica un andamento della derivata seconda come $-42x$, quindi decrescente, se \`e lecito derivare il resto. Nota: dal solo resto $o(x^4)$ non segue in generale la monotonia locale di $f''$; il quesito sottintende regolarit\`a aggiuntiva.

% Pag. originale 514/521
\iffalse Pag. originale 514/521 \fi
\begin{quiz}
Il valore dell'integrale $\displaystyle\int_0^4\frac{1}{1+\sqrt{x}}\,dx$ \`e
\begin{enumerate}
\item $0$
\item $2-4\log3$
\item $3$
\item $4-2\log3$
\item $4-2\log3+c$
\end{enumerate}
\end{quiz}
% SOLUZIONE: D
% SEZIONE_AM1: 5.4
% SPIEGAZIONE: Con $t=\sqrt{x}$, $dx=2t\,dt$, dunque $\int_0^2 2t/(1+t)\,dt=[2t-2\log(1+t)]_0^2=4-2\log3$.

% Pag. originale 514/521
\iffalse Pag. originale 514/521 \fi
\begin{quiz}
Supponendo che il teorema di integrazione per sostituzione sia applicabile, tramite la sostituzione $x=g(t)$ si ottiene
\begin{enumerate}
\item $\int f(x)\,dx=\int f(g^{-1}(x))\,dx$
\item $\int f(x)\,dx=\int f(g(t))\,dt$
\item $\int f(x)\,dx=\int f(g(t))g'(t)\,dt$
\item $\int f(x)\,dx=\int f(t)g'(t)\,dt$
\item $\int f(x)\,dx=\int f(g(t))g'(t)\,dx$
\end{enumerate}
\end{quiz}
% SOLUZIONE: C
% SEZIONE_AM1: 5.4
% SPIEGAZIONE: Da $x=g(t)$ segue $dx=g'(t)\,dt$, quindi l'integrando trasformato \`e $f(g(t))g'(t)$.

% Pag. originale 515/521
\iffalse Pag. originale 515/521 \fi
\begin{quiz}
L'integrale improprio $\displaystyle\int_0^{+\infty}\frac{\arctan x}{x^\alpha}\,dx$
\begin{enumerate}
\item converge $\forall\alpha\in\mathbb R$
\item converge $\forall\alpha<1$
\item converge $\forall\alpha\in(1,2)$
\item converge $\forall\alpha>1$
\item diverge $\forall\alpha\in\mathbb R$
\end{enumerate}
\end{quiz}
% SOLUZIONE: C
% SEZIONE_AM1: 5.5
% SPIEGAZIONE: Per $x\to0^+$, $\arctan x\sim x$, quindi serve $\alpha<2$; per $x\to+\infty$, $\arctan x\to\pi/2$, quindi serve $\alpha>1$.

% Pag. originale 515/521
\iffalse Pag. originale 515/521 \fi
\begin{quiz}
L'equazione differenziale $y''+y'=0$
\begin{enumerate}
\item ha infinite soluzioni costanti
\item ha una ed una sola soluzione costante
\item non ha soluzioni costanti
\item ha $y(x)=e^x$ come integrale particolare
\item ha $y(x;c_1,c_2)=c_1e^{-x}+c_2e^x$ come integrale generale
\end{enumerate}
\end{quiz}
% SOLUZIONE: A
% SEZIONE_AM1: 6.1
% SPIEGAZIONE: L'equazione caratteristica \`e $r(r+1)=0$; la soluzione generale \`e $y=C_1+C_2e^{-x}$, con infinite soluzioni costanti scegliendo $C_2=0$.

% Pag. originale 517/521
\iffalse Pag. originale 517/521 \fi
\begin{quiz}
Siano $A,B\subset\mathbb R$. Se $A\subset B$, allora
\begin{enumerate}
\item $\inf A>\inf B$
\item $\inf A\le\inf B$
\item $\inf A<\inf B$
\item $\inf A=\inf B$
\item $\inf A\ge\inf B$
\end{enumerate}
\end{quiz}
% SOLUZIONE: E
% SEZIONE_AM1: 1.2
% SPIEGAZIONE: Il sottoinsieme $A$ non pu\`o contenere valori inferiori all'estremo inferiore di $B$: $\inf A\ge\inf B$, quando gli estremi sono definiti.

% Pag. originale 517/521
\iffalse Pag. originale 517/521 \fi
\begin{quiz}
La funzione $f(x)=\cos x+e^x$
\begin{enumerate}
\item \`e pari
\item ha infiniti zeri
\item \`e iniettiva su $\mathbb R$
\item \`e monotona su $\mathbb R$
\item \`e periodica
\end{enumerate}
\end{quiz}
% SOLUZIONE: B
% SEZIONE_AM1: 1.4
% SPIEGAZIONE: Per $x\to-\infty$, $e^x\to0$ e $\cos x$ oscilla fra $-1$ e $1$: per continuit\`a si trovano infiniti zeri distinti. Non \`e monotona n'e periodica.

% Pag. originale 517/521
\iffalse Pag. originale 517/521 \fi
\begin{quiz}
Un polinomio a coefficienti reali che ha tra le proprie radici i numeri $0$ e $2+i$
\begin{enumerate}
\item ha grado $2$
\item ha grado maggiore o uguale a $3$
\item ha grado strettamente minore di $3$
\item ha grado $4$
\item ha grado pari
\end{enumerate}
\end{quiz}
% SOLUZIONE: B
% SEZIONE_AM1: 1.5
% SPIEGAZIONE: Per coefficienti reali anche $2-i$ \`e radice. Le radici distinte $0$, $2+i$, $2-i$ impongono grado almeno $3$.

% Pag. originale 517/521
\iffalse Pag. originale 517/521 \fi
\begin{quiz}
Sia $f:\mathbb R\to\mathbb R$ una funzione tale che $\lim_{x\to2}f(x)=1$. Allora
\begin{enumerate}
\item $f(2)=1$
\item $\forall\delta>0$, se $|x-2|<\delta$ allora $f(x)>0$
\item $\exists\delta>0$ tale che se $|x-2|<\delta$ allora $|f(x)-1|<\delta$
\item $\forall\delta>0$, se $0<|x-2|<\delta$ allora $1<f(x)<2$
\item $\exists\delta>0$ tale che se $0<|x-2|<\delta$ allora $f(x)>0$
\end{enumerate}
\end{quiz}
% SOLUZIONE: E
% SEZIONE_AM1: 2.2
% SPIEGAZIONE: Poich'e il limite \`e $1>0$, esiste un intorno bucato di $2$ in cui $f(x)>0$. Nessuna condizione \`e imposta al valore $f(2)$.

% Pag. originale 517/521
\iffalse Pag. originale 517/521 \fi
\begin{quiz}
Il limite $\displaystyle\lim_{x\to0}\frac{9^x-1}{\sqrt[4]{1+8x}-1}$ vale
\begin{enumerate}
\item $3$
\item $2$
\item $\log3$
\item $\log2$
\item $9/8$
\end{enumerate}
\end{quiz}
% SOLUZIONE: C
% SEZIONE_AM1: 2.2
% SPIEGAZIONE: $9^x-1\sim x\log9$ e $(1+8x)^{1/4}-1\sim2x$: il rapporto tende a $\log9/2=\log3$.

% Pag. originale 517/521
\iffalse Pag. originale 517/521 \fi
\begin{quiz}
Il limite $\displaystyle\lim_{x\to+\infty}(3x+2)\sin^2x$
\begin{enumerate}
\item vale $+\infty$
\item vale $3$
\item non esiste
\item vale $2$
\item vale $0$
\end{enumerate}
\end{quiz}
% SOLUZIONE: C
% SEZIONE_AM1: 2.2
% SPIEGAZIONE: Lungo $x=k\pi$ il prodotto \`e $0$, mentre lungo $x=\pi/2+k\pi$ tende a $+\infty$. Il limite non esiste.

% Pag. originale 518/521
\iffalse Pag. originale 518/521 \fi
\begin{quiz}
Se le funzioni $f,g:\mathbb R\to\mathbb R$ hanno rispettivamente le rette $y=x+1$, $y=2x$ come asintoti a $+\infty$, allora
\begin{enumerate}
\item $\lim_{x\to+\infty}(f(x)-g(x))=-1$
\item $\lim_{x\to+\infty}(f(x)-g(x))=-\infty$
\item $\lim_{x\to+\infty}\frac{f(x)-g(x)}{x}=0$
\item $\lim_{x\to+\infty}\frac{f(x)-g(x)}{x}=+\infty$
\item $\lim_{x\to+\infty}(f(x)+g(x))=-\infty$
\end{enumerate}
\end{quiz}
% SOLUZIONE: B
% SEZIONE_AM1: 2.2
% SPIEGAZIONE: $f(x)=x+1+o(1)$ e $g(x)=2x+o(1)$; pertanto $f-g=-x+1+o(1)\to-\infty$.

% Pag. originale 518/521
\iffalse Pag. originale 518/521 \fi
\begin{quiz}
Sia $a_n=\dfrac{3n+(-1)^n}{n+(-2)^n}$. Allora
\begin{enumerate}
\item $\lim_{n\to\infty}a_n=0$
\item $a_n$ non ha limite per $n\to\infty$
\item $\lim_{n\to\infty}a_n=3$
\item $\lim_{n\to\infty}a_n=+\infty$
\item $a_n$ \`e monotona
\end{enumerate}
\end{quiz}
% SOLUZIONE: A
% SEZIONE_AM1: 2.2
% SPIEGAZIONE: Poich'e $|(-2)^n|=2^n$ domina $n$ e il numeratore \`e $O(n)$, si ha $a_n\to0$ da entrambi i sottogruppi.

% Pag. originale 518/521
\iffalse Pag. originale 518/521 \fi
\begin{quiz}
Sia $f(x)=2^{x\sin x}$. Allora
\begin{enumerate}
\item $f'(x)=2^{x\sin x}(\sin x+x\cos x)$
\item $f'(x)=2^{x\sin x}$
\item $f'(x)=2^{x\sin x}x\cos x$
\item $f'(x)=2^{x\sin x}(\sin x+x\cos x)\log2$
\item $f'(x)=2^{x\sin x}\log2$
\end{enumerate}
\end{quiz}
% SOLUZIONE: D
% SEZIONE_AM1: 4.1
% SPIEGAZIONE: Derivando $a^{u(x)}$ si ottiene $a^{u(x)}\log a\,u'(x)$, con $u'=\sin x+x\cos x$.

% Pag. originale 518/521
\iffalse Pag. originale 518/521 \fi
\begin{quiz}
Se $f(x_0)=0$ ed $x_0$ \`e un punto angoloso per $f$, allora $g(x)=(x-x_0)f(x)$
\begin{enumerate}
\item \`e derivabile in $x_0$
\item ha un punto di cuspide in $x_0$
\item ha un punto angoloso in $x_0$
\item non \`e derivabile in $x_0$
\item \`e discontinua in $x_0$
\end{enumerate}
\end{quiz}
% SOLUZIONE: A
% SEZIONE_AM1: 4.1
% SPIEGAZIONE: Il rapporto incrementale di $g$ in $x_0$ \`e $f(x)$, che tende a $f(x_0)=0$ perch'e un punto angoloso implica continuit\`a. Quindi $g'(x_0)=0$.

% Pag. originale 518/521
\iffalse Pag. originale 518/521 \fi
\begin{quiz}
Sia $f:\mathbb R\to\mathbb R$ una funzione derivabile ed invertibile. Se $f(2)=1$ e $f'(2)=3$ allora
\begin{enumerate}
\item $(f^{-1})'(2)=3$
\item $(f^{-1})'(3)=1$
\item $(f^{-1})'(1)=1/3$
\item $(f^{-1})'(1)=1/2$
\item $(f^{-1})'(3)=2$
\end{enumerate}
\end{quiz}
% SOLUZIONE: C
% SEZIONE_AM1: 4.1
% SPIEGAZIONE: La formula per la derivata dell'inversa d\`a $(f^{-1})'(f(2))=1/f'(2)=1/3$.

% Pag. originale 518/521
\iffalse Pag. originale 518/521 \fi
\begin{quiz}
Sia $f:[0,4]\to\mathbb R$ definita da $f(x)=\sqrt{x}$ se $x\ne1$, e $f(x)=\alpha$ se $x=1$, con $\alpha$ parametro reale. Allora $f$ assume tutti i valori compresi tra $0$ e $2$ se e solo se
\begin{enumerate}
\item $\alpha=0$
\item $\alpha=1$
\item $\alpha=2$
\item $0\le\alpha\le4$
\item $0\le\alpha\le2$
\end{enumerate}
\end{quiz}
% SOLUZIONE: B
% SEZIONE_AM1: 1.4
% SPIEGAZIONE: Il valore $1$ \`e assunto da $\sqrt{x}$ soltanto in $x=1$, dove la funzione vale invece $\alpha$. Per coprire $[0,2]$ occorre e basta $\alpha=1$.

% Pag. originale 519/521
\iffalse Pag. originale 519/521 \fi
\begin{quiz}
Sia $f:\mathbb R\to\mathbb R$ continua su $\mathbb R$ e decrescente in $(-\infty,0)$ e in $(0,+\infty)$. Allora
\begin{enumerate}
\item $f(\mathbb R)$ non \`e un intervallo
\item $f(\mathbb R)=\mathbb R$
\item $f$ \`e decrescente su tutto $\mathbb R$
\item $\sup_{\mathbb R}f=+\infty$
\item $f$ \`e limitata su $\mathbb R$
\end{enumerate}
\end{quiz}
% SOLUZIONE: C
% SEZIONE_AM1: 4.3
% SPIEGAZIONE: La continuit\`a in $0$ permette di raccordare i due tratti decrescenti senza salti: per $x<0<y$, $f(x)\ge f(0)\ge f(y)$.

% Pag. originale 519/521
\iffalse Pag. originale 519/521 \fi
\begin{quiz}
Sia $f:\mathbb R\to\mathbb R$ una funzione derivabile tale che $f(-1)=0$ e $f(6)=3$. Allora nell'intervallo $(-1,6)$ la derivata $f'(x)$
\begin{enumerate}
\item assume il valore $7/3$ in almeno un punto
\item si annulla in almeno un punto
\item assume il valore $7/3$ in infiniti punti
\item assume il valore $3/7$ in almeno un punto
\item non \`e mai nulla
\end{enumerate}
\end{quiz}
% SOLUZIONE: D
% SEZIONE_AM1: 4.2
% SPIEGAZIONE: Il teorema di Lagrange garantisce $f'(c)=[f(6)-f(-1)]/[6-(-1)]=3/7$ per almeno un $c\in(-1,6)$.

% Pag. originale 519/521
\iffalse Pag. originale 519/521 \fi
\begin{quiz}
Per $x\to0$ si ha
\begin{enumerate}
\item $\dfrac1{x^4+x^2}\sim\dfrac1{x^4}$
\item $\dfrac1{x^4}+\dfrac1{x^2}\sim\dfrac1{x^2}$
\item $\dfrac1{x^4}-\dfrac1{x^2}\sim-\dfrac1{x^2}$
\item $x^4+x^2\sim x^2$
\item $\dfrac1{x^4-x^2}\sim\dfrac1{x^4}$
\end{enumerate}
\end{quiz}
% SOLUZIONE: D
% SEZIONE_AM1: 3.4
% SPIEGAZIONE: $x^4+x^2=x^2(1+x^2)$ e il rapporto con $x^2$ tende a $1$.

% Pag. originale 519/521
\iffalse Pag. originale 519/521 \fi
\begin{quiz}
Il limite $\displaystyle\lim_{x\to0}\frac{e^{\sqrt[3]{8x}}-2\sqrt[3]{x}-1}{\sin\sqrt[7]{x}}$ vale
\begin{enumerate}
\item $1$
\item $+\infty$
\item $0$
\item $-1$
\item $e$
\end{enumerate}
\end{quiz}
% SOLUZIONE: C
% SEZIONE_AM1: 3.4
% SPIEGAZIONE: $\sqrt[3]{8x}=2\sqrt[3]{x}$; lo sviluppo dell'esponenziale cancella il termine lineare, lasciando $O(|x|^{2/3})$. Il denominatore \`e asintotico a $\sqrt[7]{x}$, dunque il rapporto tende a $0$.

% Pag. originale 519/521
\iffalse Pag. originale 519/521 \fi
\begin{quiz}
Lo sviluppo di Mac Laurin di ordine $3$ di $f(x)=\dfrac{2x}{1+x}-\dfrac1{x-1}$ \`e
\begin{enumerate}
\item $f(x)=1-3x-x^2+3x^3+o(x^3)$
\item $f(x)=1+3x-x^2-3x^3+o(x^3)$
\item $f(x)=1+3x-x^2+3x^3+o(x^3)$
\item $f(x)=1-3x-x^2-3x^3+o(x^3)$
\item $f(x)=1+3x+x^2+3x^3+o(x^3)$
\end{enumerate}
\end{quiz}
% SOLUZIONE: C
% SEZIONE_AM1: 4.2
% SPIEGAZIONE: $2x/(1+x)=2x-2x^2+2x^3+o(x^3)$ e $-1/(x-1)=1+x+x^2+x^3+o(x^3)$; sommando si ottiene C.

% Pag. originale 519/521
\iffalse Pag. originale 519/521 \fi
\begin{quiz}
Sia $F(x)$ la primitiva di $f(x)=\dfrac{e^{\sqrt{x+2}}}{\sqrt{x+2}}$ che si annulla in $x=-1$. Allora
\begin{enumerate}
\item $F(x)=2\sqrt{x+2}e^{\sqrt{x+2}}-2e$
\item $F(x)=\sqrt{x+2}e^{\sqrt{x+2}}-e$
\item $F(x)=e^{\sqrt{x+2}}-e$
\item $F(x)=2e^{\sqrt{x+2}}-2e$
\item $F(x)=2e^{\sqrt{x+2}}$
\end{enumerate}
\end{quiz}
% SOLUZIONE: D
% SEZIONE_AM1: 5.4
% SPIEGAZIONE: Con $t=\sqrt{x+2}$, $dx=2t\,dt$, e $\int e^t/t\,2t\,dt=2e^t+C$. Da $F(-1)=0$ segue $C=-2e$.

% Pag. originale 520/521
\iffalse Pag. originale 520/521 \fi
\begin{quiz}
L'integrale improprio $\displaystyle\int_0^{+\infty}\frac{\sqrt{x}}{x+x^2}\,dx$
\begin{enumerate}
\item converge
\item diverge positivamente
\item diverge negativamente
\item \`e indeterminato
\item vale $0$
\end{enumerate}
\end{quiz}
% SOLUZIONE: A
% SEZIONE_AM1: 5.5
% SPIEGAZIONE: L'integranda \`e $1/[\sqrt{x}(1+x)]$: vicino a $0$ \`e dell'ordine di $x^{-1/2}$, all'infinito di $x^{-3/2}$; entrambi gli integrali convergono.

% Pag. originale 520/521
\iffalse Pag. originale 520/521 \fi
\begin{quiz}
L'equazione differenziale $x''-5x'+6x=e^{2t}$ ha una soluzione particolare $x_p(t)$ della forma
\begin{enumerate}
\item $x_p(t)=Cte^{3t}$ con $C\in\mathbb R$
\item $x_p(t)=Ct^2e^{3t}$ con $C\in\mathbb R$
\item $x_p(t)=Cte^{2t}$ con $C\in\mathbb R$
\item $x_p(t)=Ct^2e^{2t}$ con $C\in\mathbb R$
\item $x_p(t)=C$ con $C\in\mathbb R$
\end{enumerate}
\end{quiz}
% SOLUZIONE: C
% SEZIONE_AM1: 6.1
% SPIEGAZIONE: L'equazione caratteristica \`e $(r-2)(r-3)=0$: il forzamento $e^{2t}$ \`e risonante con una radice semplice, quindi si cerca $x_p=Cte^{2t}$; $C=-1$.
Risposta completata
